% ============================================================
% section8_angular.tex - Раздел 8: Теоретическая модель: программная реализация
% ============================================================
% !TeX program = xelatex
% Compartibility check: xelatex --shell-escape main.tex

\section{Теоретическая модель: программная реализация}

В предыдущем разделе мы разобрали теоретические основы работы проволочных поляризаторов и матричный формализм для системы из двух решеток. Теперь пришло время перевести эту электродинамическую модель на язык Python. Мы напишем модуль, реализующий расчеты по статьям Blanco et al. (1986) и Kaltenecker et al. (2019).

\subsection{Реализация формул Бланко: от теории к коду}

Для расчёта теоретического пропускания мы создадим модуль \mycode{theoretical.py} в папке \mycode{src/}. В нём нам потребуется вычислить функции $A_1$, $A_2$, $f_a$, $f_b$, $f_c$ и $f_d$, которые, в свою очередь, зависят от суммы бесконечных рядов. На практике достаточно ограничиться первыми $N=15$ членами ряда, чтобы получить очень высокую точность.

Давай сопоставим ключевые формулы из статьи Blanco с их реализацией на Python:

\textbf{1. Функция $C_m$} (Формула~(\ref{eq:Cm}) из раздела~\ref{sec:theory}):
\begin{equation*}
C_m = \left[ m^2 - \left(\frac{p}{\lambda}\right)^2 \right]^{1/2}
\end{equation*}
Поскольку для высоких частот (когда шаг решетки $p$ превышает длину волны $\lambda$) подкоренное выражение может стать отрицательным, $C_m$ становится комплексным числом. Кроме того, при точном резонансе возникает деление на ноль. Поэтому в коде мы явно добавляем бесконечно малый комплексный регулятор \mycode{1j * 1e-9}:
\begin{minted}{python}
def compute_C(m: int, p_over_lambda: float) -> complex:
    val = m**2 - p_over_lambda**2
    return np.sqrt(val + 1j * 1e-9)
\end{minted}

\textbf{2. Параметр $A_1$} (Формула~(\ref{eq:A1}) из раздела~\ref{sec:theory}):
В статье формула выглядит так:
\begin{equation*}
A_1 = 1 + \frac{1}{2}\left(\frac{\pi d}{\lambda}\right)^2 \left( \ln\frac{p}{\pi d} + \frac{3}{4} \right) + \frac{1}{2}\left(\frac{\pi d}{\lambda}\right)^2 \sum_{m=1}^\infty \left( \frac{1}{C_m} - \frac{1}{m} \right)
\end{equation*}
В коде мы разбиваем её на три слагаемых (\mycode{term1}, \mycode{term2}, \mycode{term3}) и используем защищенный логарифм \mycode{safe_log} для предотвращения математических ошибок (например, деления на ноль или логарифма от отрицательного числа), которые могут возникнуть из-за ошибок округления:
\begin{minted}{python}
def safe_log(x: float) -> float:
    return np.log(max(x, EPS))
\end{minted}

\textbf{3. Параметр $A_2$} (Формула~(\ref{eq:A2}) из раздела~\ref{sec:theory}):
По аналогии вычисляется и второй параметр:
\begin{equation*}
A_2 = 1 + \frac{1}{2}\left(\frac{\pi d}{\lambda}\right)^2 \left( \frac{11}{4} - \ln\frac{p}{\pi d} \right) + \frac{1}{24}\left(\frac{\pi d}{\lambda}\right)^2 - \left(\frac{\pi d}{\lambda}\right)^2 \sum_{m=1}^\infty \left( m - \frac{1}{2m}\left(\frac{p}{\lambda}\right)^2 - C_m \right)
\end{equation*}
Как и в предыдущем пункте, программная реализация разбивает это длинное выражение на отдельные слагаемые.

\textbf{4. Эквивалентные проводимости $f_a$ и $f_b$} (Формулы~(\ref{eq:fa}) и (\ref{eq:fb}) из раздела~\ref{sec:theory}):
Эти функции определяют мнимые части импедансов:
\begin{align*}
f_a(p, d, \lambda) &= \frac{p}{2\lambda}\left(\frac{\pi d}{p}\right)^2 \frac{1}{A_2} \\
f_b(p, d, \lambda) &= \frac{2\lambda}{p}\left(\frac{p}{\pi d}\right)^2 A_1 - \frac{p}{4\lambda}\left(\frac{\pi d}{p}\right)^2 \frac{1}{A_2}
\end{align*}
\note{Обрати внимание на опечатку в оригинальной статье Blanco (1986): там формула для $f_b$ ошибочно пронумерована как (19), хотя должна быть (9). В более свежей работе Kaltenecker (2019) эта формула приведена корректно в системе уравнений (5).}

\subsection{Вычисление амплитудных коэффициентов пропускания}

Следующим шагом является расчет комплексных амплитудных коэффициентов $t_\perp$ и $t_\parallel$ по формулам~(\ref{eq:t_perp}) и (\ref{eq:t_par}) из Blanco:

\begin{minted}{python}
def compute_t_perp(p_over_lambda: float, d_over_p: float, N: int = 15) -> complex:
    fa = compute_fa(p_over_lambda, d_over_p, N)
    fb = compute_fb(p_over_lambda, d_over_p, N)
    
    Za = -1j / fa if abs(fa) > EPS else -1e9j
    Zb = -1j / fb if abs(fb) > EPS else -1e9j
    
    # Z1/Z0 (ур. 4) и Z2/Z0 (ур. 5)
    num1 = Za**2 + Za*Zb + (Za**2)*Zb
    den1 = 2.0*Za + Za**2 + Zb + Za*Zb
    Z1 = num1 / den1 if abs(den1) > EPS else 0j
    
    Z2 = Za / (1.0 + Za) if abs(1.0 + Za) > EPS else 1.0 + 0j
    
    # T_perp (ур. 3)
    num_t = 2.0 * Z1 * Z2
    den_t = (1.0 + Z1) * (Zb + Z2)
    return num_t / den_t if abs(den_t) > EPS else 0j
\end{minted}

\subsection{Матрица Джонса для двух поляризаторов}

Система из двух поляризаторов, где первый повернут на угол $\theta$, описывается матричным уравнением~(\ref{eq:jones_system}). В Python матричные операции реализуются через оператор \mycode{@} (матричное умножение) из библиотеки NumPy:

\begin{minted}{python}
def transmission_two_polarizers(theta: float, p_over_lambda: float, d_over_p: float, N: int = 15) -> float:
    t_perp = compute_t_perp(p_over_lambda, d_over_p, N)
    t_par = compute_t_par(p_over_lambda, d_over_p, N)

    c, s = np.cos(theta), np.sin(theta)

    # Матрицы вращения (R и R_inv) и матрицы поляризатора (P)
    R = np.array([[c, -s], [s, c]])
    R_inv = np.array([[c, s], [-s, c]])
    P = np.array([[t_perp, 0.0], [0.0, t_par]])

    # Результирующая матрица Джонса: M = P * R(theta) * P * R(-theta)
    M = P @ R @ P @ R_inv

    # Входное поле поляризовано вдоль оси X
    E_in = np.array([1.0, 0.0])
    E_out = M @ E_in

    # Энергетический коэффициент пропускания (мощность)
    T = np.abs(E_out[0])**2 + np.abs(E_out[1])**2
    return float(np.clip(T, 0.0, 1.0)) # Защита от выхода за пределы [0, 1]
\end{minted}

\subsection{Модифицированная модель для реального эксперимента}

Идеальная теория Бланко предполагает абсолютно ровные, идеально проводящие бесконечные проволоки. На практике решетка имеет дефекты: неровности намотки, конечное сопротивление металла, рассеяние на шероховатостях.

Чтобы приблизить теоретическую кривую к экспериментальным данным, мы вводим \textbf{модифицированную модель} в функции \mycode{transmission_db_modified}. Она содержит два подгоночных параметра:
\begin{itemize}
    \item \mycode{period_scale} --- масштабный множитель для периода. Из-за нерегулярности намотки эффективный шаг решетки $p_{\text{eff}}$ может отличаться от паспортного значения.
    \item \mycode{loss_factor} --- коэффициент частотно-зависимых потерь. Он моделирует экспоненциальное затухание, растущее с частотой, обусловленное омическими потерями и рассеянием.
\end{itemize}

\begin{minted}{python}
def transmission_db_modified(theta_deg: float, p: float, d: float, freq_THz: float, 
                             period_scale: float, loss_factor: float, N: int = 15) -> float:
    # ...
    p_eff = p * period_scale
    T_linear = transmission_two_polarizers(theta_rad, p_eff/lambda_m, d/p_eff, N)
    
    # Добавляем эмпирические потери
    T_linear_mod = T_linear * np.exp(-loss_factor * freq_THz)
    return 10 * np.log10(np.maximum(T_linear_mod, 1e-12))
\end{minted}

\subsection{Полный листинг модуля}

Для продолжения работы полный код модуля \mycode{theoretical.py} вынесен в \textbf{\textit{Приложение~\ref{app:theoretical}}}. Этот модуль станет нашим вычислительным «ядром» для расчёта теоретических кривых. Скопируй его в файл \mycode{theoretical.py} в корне вашего проекта в папке \mycode{src/}.

\subsubsection{Разбор кода: Функции модуля theoretical.py}

Для лучшего понимания работы расчетного ядра разберем ключевые параметры и функции, используемые в модуле:

\begin{syntaxdesc}
    \item[Параметр функции \mycode{p_over_lambda: float}] Относительный период решетки, равный отношению шага решетки $p$ к длине волны падающего излучения $\lambda$.
    \item[Параметр функции \mycode{d_over_p: float}] Коэффициент заполнения (filling factor), определяющий отношение диаметра проволоки $d$ к периоду решетки $p$.
    \item[Количество гармоник \mycode{N: int = 15}] Число удерживаемых членов в бесконечных рядах Бланко для обеспечения необходимой точности расчетов.
    \item[Угол поворота \mycode{theta: float}] Угол поворота первого поляризатора в радианах, определяющий взаимную ориентацию осей в матричном формализме.
    \item[Масштаб периода \mycode{period_scale: float}] Подгоночный параметр, компенсирующий отклонение эффективного шага решетки от паспортного значения.
    \item[Коэффициент затухания \mycode{loss_factor: float}] Эмпирический параметр потерь, моделирующий частотно-зависимое омическое поглощение и рассеяние.
\end{syntaxdesc}

\subsection{Итоги раздела}
Мы успешно перевели электродинамическую модель Бланко и матричный формализм Кальтенекера в надежный Python-код. Особое внимание было уделено защите от математических сингулярностей (деление на ноль, отрицательные значения под логарифмом) и добавлению эмпирических параметров (\mycode{period_scale} и \mycode{loss_factor}) для приближения теории к реальности. Наше расчетное ядро \mycode{theoretical.py} готово к работе! В следующем разделе мы объединим эти расчеты с экспериментальными данными и построим наглядную угловую зависимость.
